% ============================================================================
% Fichier  : partie2/P02_C06_cobit_objectifs.tex
% Titre    : Les 40 Objectifs de Gouvernance et de Management
% Auteur   : MINKA MI NGUIDJOÏ Thierry Emmanuel
% Version  : 1.0 -- octobre 2026
% ============================================================================
% Sources : COBIT 2019 Framework: Governance and Management Objectives et
% Introduction and Methodology (ISACA, 2018), dossier Normes. Intitulés
% traduits par nos soins ; liens objectifs/buts relevés dans les fiches
% officielles (EDM03, APO12, BAI06, BAI07, DSS02, DSS04).
% ============================================================================

\chapter{Les 40 Objectifs de Gouvernance et de Management}
\label{chap:cobit-objectifs}
\chaptermark{Les 40 objectifs de COBIT 2019}

\epigraph{%
    « Quarante objectifs ne font pas quarante chantiers.\\
    Ils font une liste dans laquelle on choisit. »%
}{Principe de mise en œuvre de COBIT}

\niveauChapitre{Fondamental}{%
    Maîtrise du Chapitre~\ref{chap:cobit-archi} (principes, composantes et
    domaines de COBIT~2019).%
}

\begin{boxobjectifs}
\begin{itemize}
    \item Parcourir les objectifs de chaque domaine et en lire la
          structure.
    \item Rédiger la fiche processus d'un objectif (finalité, pratiques,
          entrées et sorties, RACI, indicateurs).
    \item Relier objectifs d'entreprise, objectifs d'alignement et
          objectifs COBIT (cascade des objectifs).
    \item Identifier les objectifs prioritaires d'une banque et d'une
          administration.
\end{itemize}
\end{boxobjectifs}

\bigskip

% ============================================================================
\section*{Question directrice : par où commencer~?}
% ============================================================================

Mme~EKAMBI, directrice de l'audit interne de la BANQUE SANAGA, a remis au
conseil un tableau des 40~objectifs de COBIT~2019. Devant chacun, une case
« conforme » ou « non conforme ». Trente-deux cases sont rouges. Le
président s'inquiète~: « Faut-il lancer trente-deux chantiers~? »

La réponse est non, et ce chapitre explique pourquoi. Il apprend d'abord à
\textbf{lire} un objectif (section~\ref{sec:c06-lire}), puis parcourt les
\textbf{quarante objectifs}, domaine par domaine
(sections~\ref{sec:c06-edm} à~\ref{sec:c06-dss-mea}). Il montre enfin
comment la \textbf{cascade des objectifs} relie les besoins de
l'organisation aux objectifs COBIT, et permet de choisir ceux qui comptent
(sections~\ref{sec:c06-cascade} et~\ref{sec:c06-priorites}).

\separateur

% ============================================================================
\section{Lire un Objectif COBIT}
\label{sec:c06-lire}
% ============================================================================

\subsection{L'anatomie d'une fiche}

Le guide \textit{COBIT~2019 Framework: Governance and Management
Objectives}~\cite{isaca2018cobitobjectifs} présente chaque objectif sous
une forme constante. Savoir la lire est la première compétence du
praticien.

\begin{tableaucours}{L{3.8cm}L{8.4cm}}{%
    \tableauHeader{Rubrique} &
    \tableauHeader{Contenu}
}
En-tête & Domaine, code et intitulé de l'objectif, domaine d'intérêt
(le modèle central, ou un domaine particulier) \\
Description et finalité & Ce que l'objectif couvre, et pourquoi il existe \\
Objectifs liés & Objectifs d'entreprise et objectifs d'alignement
auxquels l'objectif contribue \emph{à titre principal}, avec des exemples
d'indicateurs \\
A. Processus & Les pratiques, leurs activités (chacune rattachée à un
niveau de capacité) et des indicateurs \\
B. Structures organisationnelles & Une matrice RACI par pratique \\
C. Flux d'information & Entrées et sorties de chaque pratique, avec leur
provenance et leur destination \\
D. Personnes et compétences & Compétences requises, avec renvoi aux
référentiels de compétences \\
E. Politiques et procédures & Politiques pertinentes et leur contenu \\
F. Culture et comportement & Éléments culturels clés \\
G. Services et applications & Outils et services qui supportent
l'objectif \\
\end{tableaucours}
\vspace{-0.8em}
\begin{center}{\small\itshape Tableau~6.1. Structure d'une fiche d'objectif
COBIT~2019 (d'après ISACA, notre traduction)}\end{center}

On retrouve dans les rubriques~A à~G les sept composantes du
Chapitre~\ref{chap:cobit-archi}. Une fiche COBIT n'est donc pas la
description d'un processus~: c'est la description d'un \emph{morceau de
système de gouvernance}, avec toutes ses composantes.

\subsection{Pratiques et activités}

\begin{boxdef}
\termecle{Pratique}\index{pratique (COBIT)} (\textit{governance or
management practice})~: action clé qu'il faut accomplir pour atteindre un
objectif. Elle porte un code (par exemple DSS02.04) et un intitulé à
l'infinitif.

\medskip
\termecle{Activité}\index{activité (COBIT)} (\textit{activity})~: action
concrète qui met en œuvre une pratique. Chaque activité est rattachée à
un niveau de capacité (de~2 à~5)~: c'est ce qui permet de mesurer
la capacité d'un processus (Chapitre~\ref{chap:cobit-capacite}).
\end{boxdef}

\subsection{Rédiger une fiche processus}

Le cours présente les objectifs sous une forme condensée, la
\emph{fiche processus}, qui retient l'essentiel des rubriques officielles.
Voici la fiche complète d'un objectif central de l'exploitation.

\begin{boxobjet}{DSS02 · Gérer les demandes de service et les incidents}
\textbf{Description~:} répondre de manière rapide et efficace aux demandes
des utilisateurs et résoudre tous les types d'incidents~; rétablir le
service normal, enregistrer et satisfaire les demandes, enregistrer,
analyser, diagnostiquer, faire remonter et résoudre les incidents.

\textbf{Finalité~:} accroître la productivité et réduire les
perturbations par une résolution rapide des demandes et des incidents.

\textbf{Objectifs liés~:} EG01 (portefeuille de produits et services
compétitifs), EG08 (optimisation des processus internes)~; AG05 (services
conformes aux exigences des métiers).

\textbf{Pratiques~:} DSS02.01 définir les schémas de classification des
incidents et des demandes~; DSS02.02 enregistrer, classer et prioriser~;
DSS02.03 vérifier, approuver et satisfaire les demandes~; DSS02.04
analyser, diagnostiquer et affecter les incidents~; DSS02.05 résoudre les
incidents et rétablir le service~; DSS02.06 clore les demandes et les
incidents~; DSS02.07 suivre l'état et produire des rapports.

\textbf{Entrées et sorties (exemples)~:} entrées~: accords de niveau de
service (APO09), base des erreurs connues (DSS03)~; sorties~: tickets
classés et priorisés, rapports d'incidents, demandes de changement
(vers BAI06).

\textbf{RACI (extrait, adapté à la BANQUE SANAGA)~:} DSS02.05~: \emph{A}
au responsable de l'exploitation, \emph{R} au centre de services et aux
équipes de support, \emph{C} aux métiers concernés, \emph{I} au DSI pour
les incidents majeurs.

\textbf{Indicateurs~:} délai moyen de résolution, taux de résolution au
premier contact, nombre d'incidents réouverts, satisfaction des
utilisateurs.
\end{boxobjet}
{\small\itshape Description, finalité, objectifs liés et pratiques résumés
et traduits par nos soins à partir de COBIT~2019~\cite{isaca2018cobitobjectifs}.}

\begin{boxprocedure}
\textbf{Une fiche processus utilisable.} Pour qu'une fiche processus serve
réellement, et puisse être auditée, elle doit~:
\begin{itemize}
    \item citer l'objectif COBIT et l'édition de référence~;
    \item nommer le propriétaire du processus, au niveau de l'organisation
          (et non du référentiel)~;
    \item remplacer les rôles génériques de COBIT par les fonctions réelles
          de l'organisation dans la matrice RACI~;
    \item préciser pour chaque indicateur sa cible, sa source et sa
          fréquence~;
    \item être approuvée et datée, et revue au moins une fois par an.
\end{itemize}
\end{boxprocedure}

% ============================================================================
\section{Le Domaine de Gouvernance : EDM}
\label{sec:c06-edm}
% ============================================================================

Les cinq objectifs du domaine EDM relèvent de l'organe de gouvernance. Ils
reprennent, objectif par objectif, la définition de travail du
Chapitre~\ref{chap:gouv-entreprise}~: un cadre de gouvernance, trois
finalités et l'engagement des parties prenantes.

\begin{tableaucours}{L{1.4cm}L{4.6cm}L{6.2cm}}{%
    \tableauHeader{Code} &
    \tableauHeader{Intitulé} &
    \tableauHeader{Question pour le conseil}
}
EDM01 & Assurer la définition et le maintien du cadre de gouvernance &
Qui décide de quoi, et le dispositif fonctionne-t-il~? \\
EDM02 & Assurer la réalisation des bénéfices & Nos investissements
produisent-ils la valeur annoncée~? \\
EDM03 & Assurer l'optimisation des risques & Les risques liés aux SI
restent-ils dans notre appétence~? \\
EDM04 & Assurer l'optimisation des ressources & Avons-nous les
compétences et les moyens de notre stratégie~? \\
EDM05 & Assurer l'engagement des parties prenantes & Rendons-nous compte
à ceux qui doivent l'être~? \\
\end{tableaucours}
\vspace{-0.8em}
\begin{center}{\small\itshape Tableau~6.2. Les objectifs du domaine EDM}\end{center}

Les fiches EDM01, EDM02 et EDM05 ont été présentées aux
Chapitres~\ref{chap:gouv-entreprise}, \ref{chap:alignement}
et~\ref{chap:cobit-archi}. EDM03 a pour finalité de s'assurer que les
risques d'entreprise liés aux SI ne dépassent pas l'appétence et la
tolérance au risque de l'organisation~; il sera repris avec APO12.

% ============================================================================
\section{Aligner, Planifier et Organiser : APO}
\label{sec:c06-apo}
% ============================================================================

Le domaine APO est le plus vaste. Il couvre l'organisation d'ensemble des
SI, leur stratégie et les activités de soutien.

\begin{tableaucours}{L{1.4cm}L{4.6cm}L{6.2cm}}{%
    \tableauHeader{Code} &
    \tableauHeader{Intitulé} &
    \tableauHeader{Ce qu'il couvre}
}
APO01 & Gérer le cadre de management des SI & Politiques, rôles, choix des
référentiels (Chapitre~\ref{chap:referentiels}) \\
APO02 & Gérer la stratégie & Plan stratégique des SI, feuille de route \\
APO03 & Gérer l'architecture d'entreprise & Cohérence des processus,
données, applications et technologies \\
APO04 & Gérer l'innovation & Veille, expérimentation des technologies
émergentes \\
APO05 & Gérer le portefeuille & Arbitrage des investissements
(Chapitre~\ref{chap:alignement}) \\
APO06 & Gérer le budget et les coûts & Budget des SI, refacturation,
suivi des coûts \\
APO07 & Gérer les ressources humaines & Effectifs, compétences,
formation, dépendance à des personnes clés \\
APO08 & Gérer les relations & Relation entre la DSI et les métiers \\
APO09 & Gérer les accords de service & Catalogue de services, accords de
niveau de service \\
APO10 & Gérer les fournisseurs & Sélection, contrats, supervision des
prestataires \\
APO11 & Gérer la qualité & Système de management de la qualité des SI \\
APO12 & Gérer le risque & Identification, évaluation et traitement des
risques liés aux SI \\
APO13 & Gérer la sécurité & Système de management de la sécurité de
l'information \\
APO14 & Gérer les données & Gouvernance et qualité des données
(Chapitre~\ref{chap:donnees-ia}) \\
\end{tableaucours}
\vspace{-0.8em}
\begin{center}{\small\itshape Tableau~6.3. Les objectifs du domaine APO
(intitulés traduits par nos soins)}\end{center}

\begin{boxalert}
\textbf{Violation typique.} L'ANPP (opération MEFOU) confie l'objectif
APO10 à la direction des achats, qui le réduit à la passation des
marchés. Une fois le contrat signé, plus personne ne supervise les
prestataires.

\smallskip
\textbf{Règle~:} gérer un fournisseur ne s'arrête pas à la signature.
APO10 couvre toute la relation, jusqu'à la supervision de la
performance et des risques du prestataire, que l'ancien DSI de TSB a été
sanctionné pour avoir négligée (Chapitre~\ref{chap:defaillances}).
\end{boxalert}

% ============================================================================
\section{Bâtir, Acquérir et Implémenter : BAI}
\label{sec:c06-bai}
% ============================================================================

Le domaine BAI couvre tout ce qui fait \emph{changer} le système
d'information~: projets, acquisitions, mises en production. C'est
le domaine où se sont produites toutes les défaillances du
Chapitre~\ref{chap:defaillances}.

\begin{tableaucours}{L{1.4cm}L{4.6cm}L{6.2cm}}{%
    \tableauHeader{Code} &
    \tableauHeader{Intitulé} &
    \tableauHeader{Ce qu'il couvre}
}
BAI01 & Gérer les programmes & Pilotage des programmes et de leurs
bénéfices \\
BAI02 & Gérer la définition des exigences & Recueil et validation des
besoins des métiers \\
BAI03 & Gérer l'identification et la construction des solutions &
Conception, développement, acquisition \\
BAI04 & Gérer la disponibilité et la capacité & Dimensionnement et
performance (Chapitre~\ref{chap:itil-niveaux}) \\
BAI05 & Gérer le changement organisationnel & Accompagnement des
utilisateurs et des métiers \\
BAI06 & Gérer les changements informatiques & Évaluation, autorisation et
suivi des changements (Chapitre~\ref{chap:itil-changements}) \\
BAI07 & Gérer l'acceptation et la transition des changements & Recette,
mise en production, retour arrière \\
BAI08 & Gérer la connaissance & Capitalisation et partage des
connaissances \\
BAI09 & Gérer les actifs & Inventaire et cycle de vie des actifs \\
BAI10 & Gérer la configuration & Base de configuration et ses relations \\
BAI11 & Gérer les projets & Conduite des projets \\
\end{tableaucours}
\vspace{-0.8em}
\begin{center}{\small\itshape Tableau~6.4. Les objectifs du domaine BAI
(intitulés traduits par nos soins)}\end{center}

% ============================================================================
\section{Délivrer, Servir, Soutenir et Surveiller : DSS et MEA}
\label{sec:c06-dss-mea}
% ============================================================================

Le domaine DSS couvre l'exploitation quotidienne des services, sécurité
comprise. Le domaine MEA ferme la boucle~: il vérifie la performance et
la conformité, et fournit l'assurance.

\begin{tableaucours}{L{1.4cm}L{4.6cm}L{6.2cm}}{%
    \tableauHeader{Code} &
    \tableauHeader{Intitulé} &
    \tableauHeader{Ce qu'il couvre}
}
DSS01 & Gérer l'exploitation & Exploitation courante, ordonnancement,
environnement physique \\
DSS02 & Gérer les demandes de service et les incidents & Centre de
services, incidents (Chapitre~\ref{chap:itil-incidents}) \\
DSS03 & Gérer les problèmes & Causes profondes, erreurs connues \\
DSS04 & Gérer la continuité & Plans de continuité, tests, reprise \\
DSS05 & Gérer les services de sécurité & Protection contre les logiciels
malveillants, accès, surveillance \\
DSS06 & Gérer les contrôles des processus métier & Contrôles applicatifs
intégrés aux processus \\
MEA01 & Gérer la surveillance de la performance et de la conformité &
Indicateurs, tableaux de bord \\
MEA02 & Gérer le système de contrôle interne & Autoévaluation, suivi des
déficiences \\
MEA03 & Gérer la conformité aux exigences externes & Lois, règlements,
contrats \\
MEA04 & Gérer l'assurance & Missions d'assurance et d'audit
(Chapitre~\ref{chap:audit}) \\
\end{tableaucours}
\vspace{-0.8em}
\begin{center}{\small\itshape Tableau~6.5. Les objectifs des domaines DSS et
MEA (intitulés traduits par nos soins)}\end{center}

\begin{boxmetier}{Responsable des services IT}
\textbf{Les six objectifs du quotidien}
\begin{enumerate}
    \item DSS01, DSS02 et DSS03 forment le noyau de l'exploitation~:
          faire tourner, réparer, éviter que cela recommence.
    \item BAI06 et BAI07 encadrent tout ce qui modifie la production.
    \item APO09 fixe les engagements de service sur lesquels l'exploitation
          est jugée.
\end{enumerate}
Ces six objectifs correspondent aux pratiques d'ITIL~4 étudiées dans la
Partie~\ref{part:pratique}~: COBIT dit ce qui doit être maîtrisé, ITIL
comment le faire.
\end{boxmetier}

% ============================================================================
\section{La Cascade des Objectifs}
\label{sec:c06-cascade}
% ============================================================================

\subsection{Des besoins aux objectifs COBIT}

Comment choisir, parmi les quarante objectifs, ceux qui comptent pour une
organisation donnée~? COBIT~2019 propose un premier outil~: la
\termecle{cascade des objectifs}\index{cascade des objectifs}
(\textit{goals cascade})~\cite{isaca2018cobitintro}.

\noindent\begin{minipage}{\linewidth}
\begin{center}
\begin{tikzpicture}[
    etape/.style = {
        draw=CouleurBordDef, fill=CouleurDef, rounded corners=4pt,
        text width=2.5cm, minimum height=1.2cm, align=center,
        font=\footnotesize\bfseries\color{CouleurPrimaire}
    },
    fleche/.style = {-{Stealth[length=5pt]}, thick, color=CouleurSecondaire}
]
\node[etape] (a) at (0,0)    {Besoins des\\ parties prenantes};
\node[etape] (b) at (3.4,0)  {13 objectifs\\ d'entreprise};
\node[etape] (c) at (6.8,0)  {13 objectifs\\ d'alignement};
\node[etape] (d) at (10.2,0) {40 objectifs\\ COBIT};
\draw[fleche] (a) -- (b);
\draw[fleche] (b) -- (c);
\draw[fleche] (c) -- (d);
\end{tikzpicture}
\end{center}
\vspace{-0.4em}
\begin{center}{\small\itshape Figure~6.1. La cascade des objectifs de
COBIT~2019}\end{center}
\end{minipage}

\begin{boxdef}
\termecle{Objectif d'entreprise}\index{objectif d'entreprise}
(\textit{enterprise goal}, EG)~: but de l'organisation dans son ensemble,
classé selon les quatre dimensions du tableau de bord prospectif
(finances, clients, processus internes, apprentissage et croissance).

\medskip
\termecle{Objectif d'alignement}\index{objectif d'alignement}
(\textit{alignment goal}, AG)~: but que l'informatique et les technologies
doivent atteindre pour soutenir les objectifs d'entreprise. COBIT~5 les
appelait « objectifs liés aux SI »~; le nouveau terme souligne qu'ils ne
sont pas des objectifs internes de la DSI.
\end{boxdef}

\subsection{Les vingt-six objectifs intermédiaires}

\begin{tableaucours}{L{1.0cm}L{5.0cm}L{1.0cm}L{5.0cm}}{%
    \tableauHeader{Code} &
    \tableauHeader{Objectif d'entreprise} &
    \tableauHeader{Code} &
    \tableauHeader{Objectif d'alignement}
}
EG01 & Portefeuille de produits et services compétitifs & AG01 &
Conformité des SI et soutien à la conformité aux lois \\
EG02 & Risque métier maîtrisé & AG02 & Risque lié aux SI maîtrisé \\
EG03 & Conformité aux lois et règlements & AG03 & Bénéfices réalisés des
investissements et services \\
EG04 & Qualité de l'information financière & AG04 & Qualité de
l'information financière liée aux technologies \\
EG05 & Culture de service orientée client & AG05 & Services conformes aux
exigences des métiers \\
EG06 & Continuité et disponibilité des services métier & AG06 & Agilité
pour transformer les exigences en solutions \\
EG07 & Qualité de l'information de gestion & AG07 & Sécurité de
l'information, des infrastructures et des applications, et vie privée \\
EG08 & Optimisation des fonctionnalités des processus internes & AG08 &
Soutien aux processus métier par l'intégration des applications \\
EG09 & Optimisation des coûts des processus & AG09 & Programmes livrés dans
les délais, les budgets et la qualité \\
EG10 & Compétences, motivation et productivité du personnel & AG10 &
Qualité de l'information de gestion des SI \\
EG11 & Conformité aux politiques internes & AG11 & Conformité des SI aux
politiques internes \\
EG12 & Programmes de transformation numérique maîtrisés & AG12 &
Personnel compétent, compréhension mutuelle des technologies et des
métiers \\
EG13 & Innovation des produits et des processus & AG13 & Connaissances et
initiatives pour l'innovation \\
\end{tableaucours}
\vspace{-0.8em}
\begin{center}{\small\itshape Tableau~6.6. Objectifs d'entreprise et
objectifs d'alignement de COBIT~2019 (intitulés abrégés et traduits par
nos soins)}\end{center}

\subsection{Lire la cascade dans les fiches}

Chaque fiche d'objectif indique les objectifs d'entreprise et
d'alignement auxquels il contribue \emph{à titre principal}. En les
relevant, on peut parcourir la cascade dans les deux sens. Le relevé
suivant est tiré des fiches officielles~\cite{isaca2018cobitobjectifs}.

\begin{tableaucours}{L{2.6cm}L{4.4cm}L{5.2cm}}{%
    \tableauHeader{Objectif} &
    \tableauHeader{Objectifs d'entreprise} &
    \tableauHeader{Objectifs d'alignement}
}
EDM03 & EG02, EG06 & AG02, AG07 \\
APO12 & EG02, EG06 & AG02, AG07 \\
DSS04 & EG01, EG02, EG06, EG08 & AG05, AG07 \\
DSS02 & EG01, EG08 & AG05 \\
BAI06, BAI07 & EG01 & AG06 \\
\end{tableaucours}
\vspace{-0.8em}
\begin{center}{\small\itshape Tableau~6.7. Liens principaux de quelques
objectifs (relevés dans COBIT~2019)}\end{center}

\begin{boiteRemarque}{Pourquoi la gestion des changements n'est-elle pas reliée au risque~?}
Le Tableau~6.7 réserve une surprise. On s'attendrait à ce que BAI06 et
BAI07, au cœur de toutes les défaillances du Chapitre~\ref{chap:defaillances},
soient rattachés en priorité au risque (EG02, AG02). Or leurs liens
principaux sont l'agilité (AG06) et la compétitivité de l'offre (EG01)~:
COBIT voit d'abord dans la gestion des changements le moyen de changer
\emph{vite et bien}.

La leçon est double. D'abord, la cascade ne relève que les liens
\emph{principaux}~; elle ne suffit pas, à elle seule, à prioriser. Ensuite,
une organisation qui a connu une défaillance due à un changement doit
remonter BAI06 et BAI07 par un autre chemin~: le profil de risque et les
problèmes liés aux SI, qui sont des facteurs de conception
(Chapitre~\ref{chap:cobit-conception}). C'est ce que fera la BANQUE
SANAGA.
\end{boiteRemarque}

% ============================================================================
\section{Objectifs Prioritaires : une Banque, une Administration}
\label{sec:c06-priorites}
% ============================================================================

\subsection{La BANQUE SANAGA}

\begin{boxscenario}{Opération SANAGA. Choisir dans la liste}
La responsable GRC part des objectifs d'entreprise que le conseil a
déclarés prioritaires après la panne~: EG06 (continuité et disponibilité),
EG02 (risque métier maîtrisé) et EG03 (conformité, au regard de la COBAC).

\medskip
\textbf{Par la cascade.} EG06 et EG02 conduisent, par les fiches
officielles, à EDM03, APO12 et DSS04 (Tableau~6.7). EG03 conduit à
l'objectif d'alignement AG01, et de là aux objectifs de conformité, dont
MEA03.

\medskip
\textbf{Par l'histoire de la banque.} La panne impose d'ajouter BAI06 et
BAI07 (changements), APO10 (fournisseurs), MEA02 (suivi des
recommandations) et EDM01 (cadre de gouvernance).

\medskip
\textbf{Résultat.} Dix objectifs prioritaires, et non trente-deux~: EDM01,
EDM03, APO10, APO12, BAI06, BAI07, DSS02, DSS04, MEA02, MEA03. Les autres
objectifs ne sont pas abandonnés~: ils seront visés à un niveau de
capacité plus modeste (Chapitre~\ref{chap:cobit-capacite}).
\end{boxscenario}

\subsection{L'ANPP}

Pour l'ANPP (opération MEFOU), le point de départ est différent. Les
objectifs d'entreprise prioritaires sont EG12 (programmes de
transformation numérique maîtrisés), EG09 (optimisation des coûts) et
EG03 (conformité, au regard des règles de la gestion publique et de la
Chambre des comptes). L'histoire de l'agence (Chapitres~\ref{chap:gouv-entreprise}
et~\ref{chap:alignement}) désigne d'elle-même les objectifs de
gouvernance du portefeuille (EDM02, APO05), de conduite des programmes et
des projets (BAI01, BAI11) et de gestion des fournisseurs (APO10). Le
choix définitif se fera avec les facteurs de conception, au
Chapitre~\ref{chap:cobit-conception}, puis avec les tables de
correspondance complètes de COBIT~2019 dans le laboratoire GCI-LAB-03.

\begin{boxmetier}{Auditeur SI}
\textbf{Répondre au « trente-deux cases rouges »}
\begin{enumerate}
    \item Une grille « conforme / non conforme » sur les 40~objectifs n'a
          pas de sens~: COBIT ne fixe pas d'exigences, mais des objectifs à
          atteindre à un niveau de capacité choisi.
    \item Demander d'abord quels objectifs ont été retenus comme
          prioritaires, et selon quelle démarche.
    \item Évaluer ensuite la capacité des seuls objectifs prioritaires, au
          regard de la cible fixée.
    \item Présenter au conseil l'écart entre capacité constatée et capacité
          cible, objectif par objectif.
\end{enumerate}
Mme~EKAMBI reprendra son tableau sur ces bases au Chapitre~\ref{chap:cobit-capacite}.
\end{boxmetier}

% ============================================================================
\separateur
% ============================================================================

\begin{boxresume}
\textbf{Ce qu'il faut retenir de ce chapitre~:}
\begin{itemize}
    \item \textbf{Une fiche COBIT décrit un morceau de système de
          gouvernance}~: description, finalité, objectifs liés, puis les
          sept composantes (processus, RACI, flux d'information,
          compétences, politiques, culture, outils).
    \item \textbf{Pratiques et activités}~: chaque activité est rattachée à
          un niveau de capacité, ce qui permet la mesure.
    \item \textbf{Quarante objectifs en cinq domaines}~: EDM pour le
          conseil~; APO pour l'organisation et la stratégie~; BAI pour tout
          ce qui change~; DSS pour l'exploitation~; MEA pour la surveillance
          et l'assurance.
    \item \textbf{La cascade des objectifs} relie les besoins des parties
          prenantes aux 13~objectifs d'entreprise, aux 13~objectifs
          d'alignement, puis aux 40~objectifs COBIT.
    \item \textbf{La cascade ne suffit pas}~: elle ne relève que les liens
          principaux~; l'histoire et le profil de risque de l'organisation
          complètent le choix.
    \item \textbf{Choisir, c'est retenir une dizaine d'objectifs
          prioritaires}, et viser pour les autres une capacité plus
          modeste.
\end{itemize}
\end{boxresume}

\bigskip

\begin{boxexo}[title={\ding{45}\quad Exercice~6.1. Trouver l'objectif \niveaubase}]
Pour chacune des situations suivantes, indiquez l'objectif COBIT~2019
concerné (code et intitulé).
\begin{enumerate}[(a)]
    \item Personne ne sait combien de licences de logiciels la banque
          possède.
    \item Le même incident de messagerie se reproduit chaque lundi.
    \item Le budget informatique est dépassé de 30~\% sans que la direction
          financière en soit informée.
    \item Les recommandations de l'audit interne restent sans suite.
    \item Les métiers se plaignent de ne jamais être consultés sur les
          projets informatiques.
\end{enumerate}
\end{boxexo}

\begin{boxexo}[title={\ding{45}\quad Exercice~6.2. La fiche processus de DSS04 \niveaumoyen}]
À l'aide du guide \textit{COBIT~2019 Governance and Management Objectives}
(dossier \texttt{Normes}), rédigez la fiche processus de l'objectif DSS04
(gérer la continuité) pour la BANQUE SANAGA, sur le modèle de la fiche
DSS02 de la section~\ref{sec:c06-lire}.
\begin{enumerate}[(a)]
    \item Résumez sa description et sa finalité, sans recopier le texte.
    \item Relevez ses objectifs d'entreprise et d'alignement liés.
    \item Listez ses pratiques et proposez une matrice RACI avec les
          fonctions réelles de la banque.
    \item Proposez trois indicateurs, avec leur cible et leur source.
\end{enumerate}
\end{boxexo}

\begin{boxexo}[title={\ding{45}\quad Exercice~6.3. Les priorités de NOUN PAY \niveauexpert}]
NOUN PAY (Bafoussam) fait face à deux enjeux~: l'enquête de l'autorité des
données personnelles sur sa notation de crédit par IA, et sa dépendance à
un fournisseur unique d'informatique en nuage.
\begin{enumerate}[(a)]
    \item Choisissez trois objectifs d'entreprise prioritaires, et
          justifiez.
    \item Descendez la cascade jusqu'aux objectifs d'alignement, puis aux
          objectifs COBIT, en vous appuyant sur les fiches officielles.
    \item Ajoutez les objectifs que la cascade ne fait pas apparaître mais
          que la situation impose, et expliquez pourquoi.
    \item Limitez votre liste finale à huit objectifs au plus, et
          justifiez chaque exclusion importante.
\end{enumerate}
\end{boxexo}

\bigskip

\begin{boxinfo}
\textbf{Transition.} La cascade des objectifs n'est qu'un des facteurs qui
orientent le choix des objectifs prioritaires. Le
Chapitre~\ref{chap:cobit-conception} présente les onze facteurs de
conception et la démarche complète de COBIT~2019 pour concevoir un
système de gouvernance sur mesure, avec les niveaux de capacité cibles.
\end{boxinfo}
